\section{Assemble three output bytes from four sextets}
For a complete quartet, recovering bytes from sextets \(q_0,q_1,q_2,q_3\) gives
\[
\begin{aligned}
b_0&=(q_0\ll2)\mathbin{\mathrm{OR}}(q_1\gg4),\\
b_1&=((q_1\ll4)\mathbin{\mathrm{OR}}(q_2\gg2))\bmod256,\\
b_2&=((q_2\ll6)\mathbin{\mathrm{OR}}q_3)\bmod256.
\end{aligned}
\]
The byte stores keep the intended low eight bits. In the final padded quartet, use
zero for padding positions during this calculation, but publish only the output byte
count already established by validation. A pad does not authorize an extra zero byte.

Load and classify all four characters before writing the first output byte. This is
necessary for exact in-place decoding. Overwriting an early character while later
characters of that quartet are still unread can corrupt the values needed for its
second or third output byte.

The illustrated scalar loop retains four sextets in sixteen private frame bytes.
Its ordinary saved registers and return address occupy offsets zero through 63; the
sextet words occupy offsets 64, 68, 72 and 76 in an 80-byte frame. Once the quartet is
loaded, the first byte and the remaining-output check are:

\Needspace{16\baselineskip}
\MARSCodeLabel{Stop publication at the validated output count.}
\begin{minted}[samepage=true]{mars-idea}
    lwu $t0, 64($sp)
    lwu $t1, 68($sp)
    lwu $t2, 72($sp)
    lwu $t3, 76($sp)
    sll $t4, $t0, 2
    srl $t5, $t1, 4
    or $t4, $t4, $t5
    sb $t4, 0($s2)
    daddiu $s6, $s6, -1
    beqz $s6, decode_done
    nop
\end{minted}

Repeat the remaining-byte check after the second store. A one-byte decoded tail stops
there before either later store; a two-byte tail stops before the third. For ordinary
complete quartets, advance source by four and destination by three. The destination
therefore never overtakes unread input during exact in-place decoding.
